% Zone „Das kennst du schon“ – aus bank/koerper/zone.jsonl (zusammenbau v0.1)
\einheitenkopf[zone][Kennst du schon]{Das kennst du schon}
\setcounter{aufgabe}{0}

%% TODO Titel: Ich-kann-Satz fehlt in der Bank (Platzhalter: Kettenname) – Nr. 1
%% TODO Anweisung: ein Satz über der ersten Teilaufgabe fehlt in der Bank – Nr. 1
\begin{aufgabe}{Räumliches Vorstellen}
\begin{teile}
\teil Ein Spielwürfel wird an den Kanten aufgeschnitten und flach ausgebreitet. Aus wie vielen Quadraten besteht das Netz? \leerfeld Quadrate
\teil Eine Mauer aus Würfeln ist $7$ Würfel lang, $3$ Würfel hoch und $1$ Würfel tief. Wie viele Quadrate siehst du von vorn, wie viele von oben? vorn \leerfeld , oben \leerfeld
\end{teile}
\end{aufgabe}

%% TODO Titel: Ich-kann-Satz fehlt in der Bank (Platzhalter: Kettenname) – Nr. 2
%% TODO Anweisung: ein Satz über der ersten Teilaufgabe fehlt in der Bank – Nr. 2
\begin{aufgabe}{Flächeninhalt von Rechteck, Dreieck, Trapez und Kreis (Grundflächen)}
\begin{teile}
\teil Rechteck $7$ cm lang, $3$ cm breit – Flächeninhalt? \leerfeld[cm²]
\teil Trapez: parallele Seiten $6$ cm und $3{,}4$ cm, Höhe $5$ cm – Flächeninhalt? \leerfeld[cm²]
\end{teile}
\end{aufgabe}

%% TODO Titel: Ich-kann-Satz fehlt in der Bank (Platzhalter: Kettenname) – Nr. 3
%% TODO Anweisung: ein Satz über der ersten Teilaufgabe fehlt in der Bank – Nr. 3
\begin{aufgabe}{Längen- und Volumeneinheiten, Kubikdezimeter als Liter, Kubikzentimeter als Milliliter, Umrechnen zwischen beiden}
\begin{teile}
\teil $4$ l – wie viel ml? \leerfeld[ml]
\teil $2{,}35$ m – wie viel cm? \leerfeld[cm]
\end{teile}
\end{aufgabe}

%% TODO Titel: Ich-kann-Satz fehlt in der Bank (Platzhalter: Kettenname) – Nr. 4
%% TODO Anweisung: ein Satz über der ersten Teilaufgabe fehlt in der Bank – Nr. 4
\begin{aufgabe}{Formel nach einer Größe umstellen (V = G · h → h = V : G)}
\begin{teile}
\teil $P = q \cdot r$ – nach $r$ umstellen?
\teil $A = g \cdot t$ mit $A = 7{,}5$ und $g = 2{,}5$ – $t$? t = \leerfeld
\end{teile}
\end{aufgabe}

%% TODO Titel: Ich-kann-Satz fehlt in der Bank (Platzhalter: Kettenname) – Nr. 5
%% TODO Anweisung: ein Satz über der ersten Teilaufgabe fehlt in der Bank – Nr. 5
\begin{aufgabe}{Multiplizieren mit Dezimalzahlen und π, Runden; Quadrieren und Wurzelziehen}
\begin{teile}
\teil $0{,}5 \cdot 18$ – wie viel? \leerfeld
\teil $\pi \cdot 6$ – auf eine Stelle gerundet? \leerfeld
\end{teile}
\end{aufgabe}

%% TODO Titel: Ich-kann-Satz fehlt in der Bank (Platzhalter: Kettenname) – Nr. 6
%% TODO Anweisung: ein Satz über der ersten Teilaufgabe fehlt in der Bank – Nr. 6
\begin{aufgabe}{Prozentwert berechnen (zehn Prozent von einem Volumen mit Komma)}
\begin{teile}
\teil $50\,\%$ von $8$ l – wie viel? \leerfeld[l]
\teil $20\,\%$ von $3{,}5$ l – wie viel? \leerfeld[l]
\end{teile}
\end{aufgabe}

%% TODO Titel: Ich-kann-Satz fehlt in der Bank (Platzhalter: Kettenname) – Nr. 7
%% TODO Anweisung: ein Satz über der ersten Teilaufgabe fehlt in der Bank – Nr. 7
\begin{aufgabe}{Flächeninhalt von Rechteck, Dreieck, Trapez und Kreis (Grundflächen) – Fehler finden}
\begin{teile}
\teil Eine runde Tischplatte hat den Durchmesser $1{,}4$ m. Tom rechnet den Flächeninhalt: \rechnung{\pi \cdot 1{,}4^2 &\approx 6{,}16 \text{ m}^2} Wo ist der Fehler?
\end{teile}
\end{aufgabe}

%% TODO Titel: Ich-kann-Satz fehlt in der Bank (Platzhalter: Kettenname) – Nr. 8
%% TODO Anweisung: ein Satz über der ersten Teilaufgabe fehlt in der Bank – Nr. 8
\begin{aufgabe}{Flächeninhalt von Rechteck, Dreieck, Trapez und Kreis (Grundflächen) – selbst rechnen}
\begin{teile}
\teil Ein runder Teppich hat den Durchmesser $2{,}2$ m – Flächeninhalt auf zwei Stellen? \leerfeld[m²]
\end{teile}
\end{aufgabe}
